\section{Discussion}
\label{sec:discussion}

\subsection{Why the two harnesses give different \mvp estimates}

Our most striking finding is the asymmetry between the two
operationalisations. In the \rag harness an ecosystem of just two
LLMs from different pretraining families preserves diversity for
twelve iterations; in the fine-tuning harness, sixteen \distilgpt
copies still collapse. The mechanism is straightforward. In the
fine-tuning ecosystem the agents' \emph{parameters} drift each
iteration, and the drift compounds: every retraining step pushes the
model toward the modal mode of the synthetic data distribution, the
eigenvalues of the data covariance shrink, and the model loses its
ability to reproduce the tails. In the \rag ecosystem the parameters
are \emph{frozen}; the only thing that can collapse is the
\emph{retrieved context}, and modest retrieval randomness combined with
the agents' fixed inductive biases keeps generations diverse.

The practical implication is that \emph{the training-time link is what
matters for collapse risk, not the inference-time loop.} A web full
of LLM-generated text whose authors do not retrain on it (\eg they
call a fixed API) is in the regime where our $N{=}2$ \rag ecosystem
is safe. A web full of LLM-generated text whose authors \emph{do}
retrain on it --- LLM-as-curator pipelines and self-improvement loops
\citep{kovac2025recursive} --- is in the Shumailov collapse regime.

\subsection{What counts as ``different''?}

Holding $N{=}4$ fixed, model-family differentiation preserves
$\mathbf{1.9\times}$ the terminal inter-agent semantic distance of the
next-best axis (\axisprompt or \axisdata). The control (\axisseed) is
only modestly worse than either of those. This is the empirically
clearest answer to the ``what counts as different'' question:
\emph{inductive bias from pretraining is the only axis that
meaningfully buys collapse-resistance}. Prompts and \rag-pool
segmentation move the conditional but keep the likelihood narrow.
Different sampling seeds shift the realisation but not the
distribution.

This sharpens \citet{hodel2026epistemic}'s data-segment framing: in a
\rag ecosystem (no further training), the data segment buys very
little diversity because all four agents converge on the same
posterior conditional on different contexts. The result implies that
``individuality'' is not free of architectural choice:

\begin{quote}
\itshape
Population diversity that defends against collapse must come from
diversity in the \emph{prior} (architecture / pretraining), not just
from diversity in the \emph{conditioning} (prompt, retrieval pool,
or random seed).
\end{quote}

\subsection{The shape of the diversity benefit in fine-tuning}

The marginal benefit of doubling $N$ is not monotonically diminishing
in our E1 data. Going $N{=}1 \to 2$ buys a $40\%$ reduction in the
terminal collapse multiplier; $N{=}2 \to 4$ buys $17\%$;
$N{=}4 \to 8$ buys $17\%$; $N{=}8 \to 16$ buys $31\%$ (computed off
\tabref{tab:e1_perp}). The non-monotonicity is interesting but
expected to be sample-noisy at our two-seed budget. What is robust is
that no $N$ tested has a perplexity slope statistically
indistinguishable from zero, and Hodel \& West's qualitative finding
that the optimal $N$ \emph{grows} with iteration count
\citep{hodel2026epistemic} is consistent with our slope reductions.

Three structural caveats explain why we see no flat curve: (i)
per-ecosystem token budget is held fixed across $N$, so larger $N$
means each agent sees fewer training tokens per generation; (ii) we
deliberately ran replace mode without real-data refresh, the harshest
\citet{shumailov2024nature} regime; (iii) all sixteen agents share one
base architecture, so there is no architecture-axis benefit. A
follow-up combining replace-mode with $N$ different pretrained bases
would test whether the architecture-axis benefit observed in E2
transfers to the parameter-updating regime.

\subsection{Limitations}

\para{Small models.} \distilgpt (82M) is the only fine-tuned backbone.
Larger models exhibit different collapse dynamics
\citep{dohmatob2024tale, hodel2026epistemic}.
\para{Replace mode only.} We deliberately avoided
\citet{gerstgrasser2024accumulating}-style accumulation to maximise the
collapse signal; combining accumulation with population diversity is
open.
\para{Short horizons.} $T{=}6$ (E1) and $T{=}12$ (E2) are short.
\citet{hodel2026epistemic}'s optimal $M$ keeps growing past $T{=}10$,
so our failure to detect collapse in $N{=}2$ \rag may reflect the
limit of our horizon, not of the ecosystem.
\para{Limited seeds.} Two seeds in E1 and three in E2/E3 -- enough
for means but not for tight error bars.
\para{API non-determinism.} OpenRouter routes to different providers;
seeds control only Python-level RNG.
\para{\rag ``data-segment'' is an analogue.} Our \axisdata partitions
the retrieval pool, not the training data; comparable in spirit to
\citet{hodel2026epistemic} but not identical.
\para{Single embedding model.} All semantic-distance metrics use one
encoder (\minilm); qualitative conclusions should be robust to
encoder choice but absolute numbers will not be.

\subsection{Broader implications}

The model-collapse literature is regularly cited as a reason to worry
about future foundation-model training on a web increasingly populated
by LLM output \citep{shumailov2024nature, dohmatob2024tale}. Our
results refine that worry into two distinguishable claims. First,
\emph{if} the LLMs producing that text are themselves retrained on
their pooled outputs, $N \le 16$ at fixed compute is not enough to
halt collapse, and our results say nothing reassuring about $N$ in the
hundreds either; the rate slows but the slope stays positive. Second,
\emph{if} the producers are a stable population of architecturally
distinct frozen models --- close to today's API-served model market ---
the substrate stays diverse with as few as $N{=}2$ distinct families.
This shifts the policy question from ``how many LLMs do we need on the
web'' to ``how many \emph{architecturally distinct} LLMs do we need,
and do we need to constrain when their outputs feed back into
training.'' The same argument cautions against monocultures even at
high $N$: a thousand fine-tuned descendants of one base model are not
a population, they are one model with extra noise.
